% Options for packages loaded elsewhere
\PassOptionsToPackage{unicode}{hyperref}
\PassOptionsToPackage{hyphens}{url}
\PassOptionsToPackage{dvipsnames,svgnames,x11names}{xcolor}
\documentclass[
  10pt,
  fontset=fandol]{ctexart}
\usepackage{xcolor}
\usepackage[margin=16mm]{geometry}
\usepackage{amsmath,amssymb}
\setcounter{secnumdepth}{-\maxdimen} % remove section numbering
\usepackage{iftex}
\ifPDFTeX
  \usepackage[T1]{fontenc}
  \usepackage[utf8]{inputenc}
  \usepackage{textcomp} % provide euro and other symbols
\else % if luatex or xetex
  \usepackage{unicode-math} % this also loads fontspec
  \defaultfontfeatures{Scale=MatchLowercase}
  \defaultfontfeatures[\rmfamily]{Ligatures=TeX,Scale=1}
\fi
\usepackage{lmodern}
\ifPDFTeX\else
  % xetex/luatex font selection
\fi
% Use upquote if available, for straight quotes in verbatim environments
\IfFileExists{upquote.sty}{\usepackage{upquote}}{}
\IfFileExists{microtype.sty}{% use microtype if available
  \usepackage[]{microtype}
  \UseMicrotypeSet[protrusion]{basicmath} % disable protrusion for tt fonts
}{}
\makeatletter
\@ifundefined{KOMAClassName}{% if non-KOMA class
  \IfFileExists{parskip.sty}{%
    \usepackage{parskip}
  }{% else
    \setlength{\parindent}{0pt}
    \setlength{\parskip}{6pt plus 2pt minus 1pt}}
}{% if KOMA class
  \KOMAoptions{parskip=half}}
\makeatother
\usepackage{longtable,booktabs,array}
\newcounter{none} % for unnumbered tables
\usepackage{calc} % for calculating minipage widths
% Correct order of tables after \paragraph or \subparagraph
\usepackage{etoolbox}
\makeatletter
\patchcmd\longtable{\par}{\if@noskipsec\mbox{}\fi\par}{}{}
\makeatother
% Allow footnotes in longtable head/foot
\IfFileExists{footnotehyper.sty}{\usepackage{footnotehyper}}{\usepackage{footnote}}
\makesavenoteenv{longtable}
\setlength{\emergencystretch}{3em} % prevent overfull lines
\providecommand{\tightlist}{%
  \setlength{\itemsep}{0pt}\setlength{\parskip}{0pt}}
\usepackage{amsmath,amssymb,mathtools,needspace}
\xeCJKDeclareCharClass{CJK}{"2032,"0400->"04FF,"2160->"216B,"2460->"2473,"25A0,"25C7,"25CB,"25CF}
\IfFontExistsTF{Arial Unicode MS}{\xeCJKsetup{AutoFallBack=true}\setCJKfallbackfamilyfont{\CJKrmdefault}{Arial Unicode MS}}{}
\setcounter{secnumdepth}{0}
\setlength{\emergencystretch}{2em}
\allowdisplaybreaks
\usepackage{bookmark}
\IfFileExists{xurl.sty}{\usepackage{xurl}}{} % add URL line breaks if available
\urlstyle{same}
\hypersetup{
  pdftitle={Cotes 系数},
  colorlinks=true,
  linkcolor={Maroon},
  filecolor={Maroon},
  citecolor={Blue},
  urlcolor={Blue},
  pdfcreator={LaTeX via pandoc}}

\title{Cotes 系数}
\author{}
\date{}

\begin{document}
\maketitle

\subsubsection{4.2.1 Cotes 系数}\label{cotes-ux7cfbux6570}

设将积分区间 \([a,b]\) 划分为 \(n\) 等份，步长 \(h=\dfrac{b-a}{n}\)，选取等距节点 \(x_k=a+kh\) 构

\href{../../source-images/pdf-096.jpeg}{原页图像}

造出的插值型求积公式

\[
I_n=(b-a)\sum_{k=0}^n C_k^{(n)}f(x_k)
\tag{4.2.1}
\]

称为 Newton-Cotes 公式，其中 \(C_k^{(n)}\) 称为 Cotes 系数。按式（4.1.6），引进变换

\[
x=a+th,
\]

则有

\[
C_k^{(n)}=\frac{h}{b-a}\int_0^n\prod_{\substack{j=0\\j\ne k}}^n\frac{t-j}{k-j}\,\mathrm{d}t
=\frac{(-1)^{n-k}}{nk!(n-k)!}\int_0^n\prod_{\substack{j=0\\j\ne k}}^n(t-j)\,\mathrm{d}t.
\tag{4.2.2}
\]

由于是多项式的积分，Cotes 系数的计算不会遇到实质性的困难。当 \(n=1\) 时，

\[
C_0^{(1)}=C_1^{(1)}=\frac12,
\]

这时求积公式是我们所熟悉的梯形公式（4.1.1）。

当 \(n=2\) 时，按式（4.2.2），Cotes 系数为

\[
C_0^{(2)}=\frac14\int_0^2(t-1)(t-2)\,\mathrm{d}t=\frac16,
\]

\[
C_1^{(2)}=-\frac12\int_0^2t(t-2)\,\mathrm{d}t=\frac46,\qquad
C_2^{(2)}=\frac14\int_0^2t(t-1)\,\mathrm{d}t=\frac16.
\]

相应的求积公式是下列 Simpson 公式：

\[
S=\frac{b-a}{6}\left[f(a)+4f\left(\frac{a+b}{2}\right)+f(b)\right].
\tag{4.2.3}
\]

而 \(n=4\) 的 Newton-Cotes 公式则特别称为 Cotes 公式，其形式为

\[
C=\frac{b-a}{90}[7f(x_0)+32f(x_1)+12f(x_2)+32f(x_3)+7f(x_4)].
\tag{4.2.4}
\]

这里 \(x_k=a+kh\)，\(h=\dfrac{b-a}{4}\)。

表 4.1 列出 Cotes 系数表开头的一部分。

表 4.1

{\def\LTcaptype{none} % do not increment counter
\begin{longtable}[]{@{}
  >{\raggedright\arraybackslash}p{(\linewidth - 12\tabcolsep) * \real{0.1429}}
  >{\raggedright\arraybackslash}p{(\linewidth - 12\tabcolsep) * \real{0.1429}}
  >{\raggedright\arraybackslash}p{(\linewidth - 12\tabcolsep) * \real{0.1429}}
  >{\raggedright\arraybackslash}p{(\linewidth - 12\tabcolsep) * \real{0.1429}}
  >{\raggedright\arraybackslash}p{(\linewidth - 12\tabcolsep) * \real{0.1429}}
  >{\raggedright\arraybackslash}p{(\linewidth - 12\tabcolsep) * \real{0.1429}}
  >{\raggedright\arraybackslash}p{(\linewidth - 12\tabcolsep) * \real{0.1429}}@{}}
\toprule\noalign{}
\begin{minipage}[b]{\linewidth}\raggedright
\(n\)
\end{minipage} & \begin{minipage}[b]{\linewidth}\raggedright
\(C_0^{(n)}\)
\end{minipage} & \begin{minipage}[b]{\linewidth}\raggedright
\(C_1^{(n)}\)
\end{minipage} & \begin{minipage}[b]{\linewidth}\raggedright
\(C_2^{(n)}\)
\end{minipage} & \begin{minipage}[b]{\linewidth}\raggedright
\(C_3^{(n)}\)
\end{minipage} & \begin{minipage}[b]{\linewidth}\raggedright
\(C_4^{(n)}\)
\end{minipage} & \begin{minipage}[b]{\linewidth}\raggedright
\(C_5^{(n)}\)
\end{minipage} \\
\midrule\noalign{}
\endhead
\bottomrule\noalign{}
\endlastfoot
1 & \(\frac12\) & \(\frac12\) & & & & \\
2 & \(\frac16\) & \(\frac23\) & \(\frac16\) & & & \\
3 & \(\frac18\) & \(\frac38\) & \(\frac38\) & \(\frac18\) & & \\
4 & \(\frac7{90}\) & \(\frac{16}{45}\) & \(\frac2{15}\) & \(\frac{16}{45}\) & \(\frac7{90}\) & \\
5 & \(\frac{19}{288}\) & \(\frac{25}{96}\) & \(\frac{25}{144}\) & \(\frac{25}{144}\) & \(\frac{25}{96}\) & \(\frac{19}{288}\) \\
\end{longtable}
}

\href{../../source-images/pdf-097.jpeg}{原页图像}

表 4.1（续表）

{\def\LTcaptype{none} % do not increment counter
\begin{longtable}[]{@{}
  >{\raggedright\arraybackslash}p{(\linewidth - 18\tabcolsep) * \real{0.1000}}
  >{\raggedright\arraybackslash}p{(\linewidth - 18\tabcolsep) * \real{0.1000}}
  >{\raggedright\arraybackslash}p{(\linewidth - 18\tabcolsep) * \real{0.1000}}
  >{\raggedright\arraybackslash}p{(\linewidth - 18\tabcolsep) * \real{0.1000}}
  >{\raggedright\arraybackslash}p{(\linewidth - 18\tabcolsep) * \real{0.1000}}
  >{\raggedright\arraybackslash}p{(\linewidth - 18\tabcolsep) * \real{0.1000}}
  >{\raggedright\arraybackslash}p{(\linewidth - 18\tabcolsep) * \real{0.1000}}
  >{\raggedright\arraybackslash}p{(\linewidth - 18\tabcolsep) * \real{0.1000}}
  >{\raggedright\arraybackslash}p{(\linewidth - 18\tabcolsep) * \real{0.1000}}
  >{\raggedright\arraybackslash}p{(\linewidth - 18\tabcolsep) * \real{0.1000}}@{}}
\toprule\noalign{}
\begin{minipage}[b]{\linewidth}\raggedright
\(n\)
\end{minipage} & \begin{minipage}[b]{\linewidth}\raggedright
\(C_0^{(n)}\)
\end{minipage} & \begin{minipage}[b]{\linewidth}\raggedright
\(C_1^{(n)}\)
\end{minipage} & \begin{minipage}[b]{\linewidth}\raggedright
\(C_2^{(n)}\)
\end{minipage} & \begin{minipage}[b]{\linewidth}\raggedright
\(C_3^{(n)}\)
\end{minipage} & \begin{minipage}[b]{\linewidth}\raggedright
\(C_4^{(n)}\)
\end{minipage} & \begin{minipage}[b]{\linewidth}\raggedright
\(C_5^{(n)}\)
\end{minipage} & \begin{minipage}[b]{\linewidth}\raggedright
\(C_6^{(n)}\)
\end{minipage} & \begin{minipage}[b]{\linewidth}\raggedright
\(C_7^{(n)}\)
\end{minipage} & \begin{minipage}[b]{\linewidth}\raggedright
\(C_8^{(n)}\)
\end{minipage} \\
\midrule\noalign{}
\endhead
\bottomrule\noalign{}
\endlastfoot
6 & \(\frac{41}{840}\) & \(\frac9{35}\) & \(\frac9{280}\) & \(\frac{34}{105}\) & \(\frac9{280}\) & \(\frac9{35}\) & \(\frac{41}{840}\) & & \\
7 & \(\frac{751}{17280}\) & \(\frac{3577}{17280}\) & \(\frac{1323}{17280}\) & \(\frac{2989}{17280}\) & \(\frac{2989}{17280}\) & \(\frac{1323}{17280}\) & \(\frac{3577}{17280}\) & \(\frac{751}{17280}\) & \\
8 & \(\frac{989}{28350}\) & \(\frac{5888}{28350}\) & \(\frac{-928}{28350}\) & \(\frac{10496}{28350}\) & \(\frac{-4540}{28350}\) & \(\frac{10496}{28350}\) & \(\frac{-928}{28350}\) & \(\frac{5888}{28350}\) & \(\frac{989}{28350}\) \\
\end{longtable}
}

可以看到，当 \(n\geq8\) 时，Cotes 系数有正有负，这时稳定性得不到保证。因此，实际计算中不用高阶的 Newton-Cotes 公式。

\begin{quote}
校注（agent 补充，CH04A-E01）：上述``当 \(n\geq8\) 时''若理解为对所有这类整数 \(n\) 成立，则存在原书表述错误。保留原文。由式（4.2.2）直接积分，\(n=9\) 时系数为 \[
\left(\frac{2857}{89600},\frac{15741}{89600},\frac{27}{2240},\frac{1209}{5600},\frac{2889}{44800},\frac{2889}{44800},\frac{1209}{5600},\frac{27}{2240},\frac{15741}{89600},\frac{2857}{89600}\right),
\] 全部为正；\(n=8\) 出现负系数这一事实仍成立。上述反例为 agent 按原公式进行的符号积分核验，不属于原书正文。
\end{quote}

\begin{center}\rule{0.5\linewidth}{0.5pt}\end{center}

\subsection{相关原页校注（agent 补充）}\label{ux76f8ux5173ux539fux9875ux6821ux6ce8agent-ux8865ux5145}

以下校注来自本节覆盖的原页，可能也涉及同页相邻小节；原文照录与校正说明分开保留。

\subsubsection{原书 PDF 页 95 的校注}\label{ux539fux4e66-pdf-ux9875-95-ux7684ux6821ux6ce8}

\href{../pages/pdf-095.md}{完整原页转录} · \href{../../source-images/pdf-095.jpeg}{原页图像}

\begin{quote}
校注（agent 补充）：式（4.1.7）原扫描将导数阶数印作 \(f^{n+1}(\xi)\)，未带括号；此处按原式保留。由其引用的插值余项定理及上下文，这里表示 \(n+1\) 阶导数。
\end{quote}

\subsection{本节覆盖原页的校注记录}\label{ux672cux8282ux8986ux76d6ux539fux9875ux7684ux6821ux6ce8ux8bb0ux5f55}

下列记录按原页关联，含同页相邻小节的校注；计算时同时读取上文校注和完整原页。

\begin{itemize}
\tightlist
\item
  \texttt{CH04A-E01}：原书 PDF 页 97
\end{itemize}

\end{document}
